(1) 直接代入定义 $Z=pv/(RT)=(1/v+B/v^2)v=1+B/v$。正维里系数表示在该截断状态方程下同温、同密度时压力高于理想气体。推导所用条件是单摩尔量、$B$ 常数，以及截断近似有效，不能延伸到 $v$ 很小时的强关联流体。
(2) 按气体对外做功为正的号约定，$W=\int_{v_1}^{v_2}p(v)dv=RT\int_{v_1}^{v_2}(v^{-1}+Bv^{-2})dv$。积分得 $W=RT\left[\ln(v_2/v_1)+B(1/v_1-1/v_2)\right]$，第一项为理想气体等温膨胀功，第二项为二阶维里修正。
(3) 由 Maxwell 关系有 $(\partial S/\partial v)_T=R(1/v+B/v^2)$，故 $\Delta S=R[\ln(v_2/v_1)+B(1/v_1-1/v_2)]$。热力学恒等式 $(\partial U/\partial v)_T=T(\partial p/\partial T)_v-p=0$，因为 $B$ 不随 $T$ 变化，因此可逆等温过程中 $\Delta U=0$、$Q=W=T\Delta S$。不能对一般温度依赖的 $B(T)$ 仍无条件使用此结论。
(4) 由状态方程求 $(\partial p/\partial v)_T=-RT/v^2-2RTB/v^3=-RT(v+2B)/v^3$。因此 $\kappa_T=-1/[v(\partial p/\partial v)_T]=v^2/[RT(v+2B)]>0$。取 $B\to0$ 时 $Z\to1$、$W\to RT\ln(v_2/v_1)$、$\Delta S\to R\ln(v_2/v_1)$、$\kappa_T\to v/(RT)=1/p$，与理想气体完全一致。
